\paragraph{A polynomial monotone Boolean circuit for the same support.}
The formal-polynomial lower bound above persists alongside a much smaller
Boolean support circuit.  For each unary or binary factor tuple, let its
weight be an input bit recording whether that tuple is allowed.  Write
$a=(a_{v,b},a_{e,b,c})$ for all these bits.  Then
\[
 \operatorname{Supp}_G(a)=
 \bigvee_{x\in\{0,1\}^V}
       \left(\bigwedge_v a_{v,x_v}\ \wedge\!
                         \bigwedge_{e=uv}a_{e,x_u,x_v}\right)
\]
has a binary monotone Boolean circuit with $O(n^3+|E_G|)$ gates,
constructible in that time from $G$.  In particular, on the bounded-degree
expander family it has size $O(n^3)$, although every monotone arithmetic
circuit for the full formal counting polynomial needs $2^{\kappa n}$
gates.  No use of E is needed for this Boolean upper bound.

Here is the full construction and proof.  First use input bits
$f_{v,b},f_{e,b,c}$ recording that a tuple is forbidden.  A forbidden
unary tuple $(v,b)$ supplies the clause $x_v\ne b$.  A forbidden edge
tuple $(e=uv,b,c)$ supplies the clause
\[
                       (x_u\ne b)\vee(x_v\ne c).
\]
Thus consistency of the allowed tuples is exactly satisfiability of
this two-CNF.  It is useful that the set of possible clauses is fixed
by $G$: the input bits indicate which clauses are present.

Make an implication graph whose $2n$ vertices are literals
$L_{v,b}$ meaning $x_v=b$, with complement
$\overline{L_{v,b}}=L_{v,1-b}$.  A unary forbidden bit supplies
the edge $L_{v,b}\to L_{v,1-b}$.  A binary forbidden bit supplies
the two edges
\[
 L_{u,b}\longrightarrow L_{v,1-c},\qquad
 L_{v,c}\longrightarrow L_{u,1-b}.
\]
Each adjacency bit is the disjunction of the forbidden-tuple bits
that supply that edge.  The total number of such occurrences is
$2n+8|E_G|$, so all adjacency bits require $O(n+|E_G|)$ binary
OR gates, allowing constants for absent edges.  Every edge
$L\to M$ has its contrapositive edge
$\overline M\to\overline L$.

Enumerate the $2n$ literals as $1,\ldots,2n$.  Compute reachability
using the following Boolean circuits:
\[
 R^{(0)}_{ij}=A_{ij}\vee[i=j],\qquad
 R^{(k)}_{ij}=R^{(k-1)}_{ij}\vee
             \bigl(R^{(k-1)}_{ik}\wedge R^{(k-1)}_{kj}\bigr)
                     \quad(1\le k\le2n).
\]
Induction on $k$ shows that $R^{(k)}_{ij}$ records whether there is
a directed path from $i$ to $j$ whose internal vertices are among
$1,\ldots,k$.  The induction uses a simple path when there is a
path, so that a visit to $k$ splits it into the two required parts.
Consequently $R^{(2n)}$ is the reachability relation.  This circuit
uses $O(n^3)$ binary AND and OR gates.  Its inputs and operations
are all monotone in the forbidden-tuple bits.

The two-CNF is unsatisfiable precisely when
\[
 U_G(f)=\bigvee_{v\in V}
 \left(R^{(2n)}_{L_{v,0},L_{v,1}}\wedge
       R^{(2n)}_{L_{v,1},L_{v,0}}\right)=1.
\]
This is exactly the ordinary two-CNF strong-component criterion of
Aspvall--Plass--Tarjan, Theorem~1, printed p.~122, cited precisely
below.  We import that classical theorem rather than reprove it.

The entire circuit for $U_G$ uses only AND, OR, constants and the
forbidden inputs.  Form its Boolean dual by interchanging AND with
OR, interchanging constants zero with one, and replacing each
forbidden input label by the corresponding allowed input label.
Induction gate by gate, using De Morgan's laws, shows that the
dual circuit computes
\[
                            1-U_G(1-a)=\operatorname{Supp}_G(a).
\]
This is an actual monotone circuit on the allowed bits; no negation
gates are introduced at its inputs or internally.  Dualization
preserves the number of gates.  This proves the claimed bound.

The classical source for the implication graph and its exact
strong-component criterion is Aspvall--Plass--Tarjan,
\emph{A linear-time algorithm for testing the truth of certain
quantified Boolean formulas}, \emph{Information Processing Letters}
8 (1979), 121--123: the graph is defined on printed p.~121, and
Theorem~1 with its complete proof is on printed p.~122.
The primary article is available at
\url{https://mathweb.ucsd.edu/~sbuss/CourseWeb/Math268_2007WS/2SAT.pdf}.
Only the ordinary existential two-CNF theorem is used here; the
quantified extension in its Theorem~2 is not used.  The cubic circuit
construction and monotone dualization are explicit above, so no
circuit bound is being imported from the article's
linear-time random-access-machine algorithm.

\paragraph{Geometric form of the support contrast.}
For each fixed allowed-tuple bit string $a$, let
\[
 B_n=\mathbb Z[x_v:v\in V]/(x_v^2-x_v:v\in V)
                   \simeq\prod_{x\in\{0,1\}^V}\mathbb Z.
\]
For a vertex or edge factor $f$, let $p_{f,a}\in B_n$ be its
Boolean interpolation polynomial, taking value $a_{f,x|I_f}$
on section $x$.  Explicitly the unary interpolation is
$a_{v,0}(1-x_v)+a_{v,1}x_v$, and the binary interpolation is
the corresponding sum of four degree-two Boolean indicator
monomials multiplied by their allowed bits.  Each $p_{f,a}$
is an idempotent.  Put $e_a=\prod_f p_{f,a}$ and
\[
                 Y_a=\operatorname{Spec}(B_n/(1-e_a)).
\]
This is the finite etale scheme over $\mathbb Z$ consisting of
exactly the allowed assignment sections.  Equivalently one may
use all equations $p_{f,a}=1$ together with the Boolean equations:
for idempotents the ideals $(1-e_a)$ and $(1-p_{f,a}:f)$ agree.
One inclusion follows by expanding $1-\prod_f p_{f,a}$ as a
telescoping sum; the other follows from
$(1-p_{f,a})e_a=0$.

For every field $K$, the scheme $Y_a\times\operatorname{Spec}K$
is nonempty if and only if $\operatorname{Supp}_G(a)=1$.
Thus the nonemptiness of this explicitly given finite etale
subscheme is computable by the displayed $O(n^3)$ monotone
Boolean circuit in its allowed-tuple parameters.  Its rank is
$P_G(a)$ as an ordinary nonnegative integer, while the full
torus character before specialization is the exponentially
costly positive arithmetic object described above.  The map
from positive formal polynomials to Boolean support is a
semiring map with idempotent addition.  It is not an evaluation
of invertible torus characters at zero in the representation
ring, and it does not retain multiplicities.  This distinction
is exactly why the arithmetic lower bound and the Boolean
upper bound coexist.
